\section{Análisis vertical}

\begin{frame}{Análisis vertical: definición}
  \begin{columns}[T,onlytextwidth]
    \column{0.46\textwidth}
    \footnotesize
    El análisis vertical estudia la composición interna de un estado financiero en un solo periodo.

    \vspace{0.5em}
    Cada partida se expresa como porcentaje de una base común, lo que permite conocer su peso relativo.
    \tiempo{0.7 minutos}
    \column{0.50\textwidth}
    \figpng[0.64\textheight]{F08_Analisis_Vertical_Base_100.png}
  \end{columns}
  \note{Definir análisis vertical.}
\end{frame}

\begin{frame}{Utilidad del análisis vertical}
  \begin{columns}[T,onlytextwidth]
    \column{0.48\textwidth}
    \begin{itemize}
      \item Medir peso de costos y gastos.
      \item Evaluar estructura de activos.
      \item Analizar deuda y patrimonio.
      \item Comparar márgenes.
    \end{itemize}
    \micro{Permite estudiar qué partes concentran recursos, costos o financiamiento dentro del estado financiero.}
    \tiempo{0.7 minutos}
    \column{0.48\textwidth}
    \figpng[0.62\textheight]{F11_Interpretacion_Vertical.png}
  \end{columns}
  \note{Explicar utilidad vertical.}
\end{frame}

\begin{frame}{Fórmula del análisis vertical}
  \begin{columns}[T,onlytextwidth]
    \column{0.48\textwidth}
    \[
      PP=\frac{\text{Partida}}{\text{Base común}}\times 100
    \]
    \begin{itemize}
      \item Resultados: ventas netas.
      \item Situación financiera: total activos.
    \end{itemize}
    \tiempo{0.8 minutos}
    \column{0.48\textwidth}
    \figpng[0.62\textheight]{F09_Formula_Analisis_Vertical.png}
  \end{columns}
  \note{Explicar fórmula vertical.}
\end{frame}

\begin{frame}{Ejemplo básico de cálculo vertical}
  \begin{columns}[T,onlytextwidth]
    \column{0.48\textwidth}
    \footnotesize
    Una empresa desea conocer qué proporción de sus ventas se destina al costo de ventas.
    \vspace{0.35em}
    \begin{tabular}{lr}
      \toprule
      Concepto & Monto \\
      \midrule
      Ventas netas & S/ 100,000.00 \\
      Costo ventas & S/ 60,000.00 \\
      Participación & 60.00 \% \\
      \bottomrule
    \end{tabular}
    \micro{De cada S/ 100.00 vendidos, S/ 60.00 se destinan al costo de ventas.}
    \tiempo{0.8 minutos}
    \column{0.48\textwidth}
    \figpng[0.58\textheight]{F08_Analisis_Vertical_Base_100.png}
  \end{columns}
  \note{Ejemplo básico vertical.}
\end{frame}

\begin{frame}{Ejemplo aplicado: estructura de costos tecnológicos}
  \begin{columns}[T,onlytextwidth]
    \column{0.43\textwidth}
    \footnotesize
    Una empresa tecnológica desea evaluar qué peso tienen sus costos de servidores y licencias dentro de sus ventas.
    \vspace{0.25em}
    \begin{itemize}
      \item Ventas: S/ 200,000.00.
      \item Servidores y licencias: S/ 70,000.00.
      \item Participación: 35.00 \%.
    \end{itemize}
    \micro{Si el porcentaje aumenta, deben revisarse arquitectura, proveedores cloud o licencias.}
    \tiempo{0.8 minutos}
    \column{0.53\textwidth}
    \figpng[0.66\textheight]{F19_Ejemplo_Vertical_Costos_Tecnologicos.png}
  \end{columns}
  \note{Ejemplo tecnológico vertical.}
\end{frame}

\begin{frame}{Procedimiento del análisis vertical}
  \begin{columns}[T,onlytextwidth]
    \column{0.48\textwidth}
    \begin{enumerate}
      \item Elegir estado financiero.
      \item Definir base 100\%.
      \item Dividir cada partida entre la base.
      \item Comparar porcentajes.
      \item Interpretar estructura.
    \end{enumerate}
    \tiempo{0.8 minutos}
    \column{0.48\textwidth}
    \figpng[0.60\textheight]{F08_Analisis_Vertical_Base_100.png}
  \end{columns}
  \note{Procedimiento vertical.}
\end{frame}

\begin{frame}{Caso Empresa ABC: cálculo vertical}
  \begin{columns}[T,onlytextwidth]
    \column{0.50\textwidth}
    \scriptsize
    \begin{tabular}{lrr}
      \toprule
      Partida & 2024 \% & 2025 \% \\
      \midrule
      Ventas & 100.00 \% & 100.00 \% \\
      Costo ventas & 60.00 \% & 61.86 \% \\
      Utilidad bruta & 40.00 \% & 38.14 \% \\
      Utilidad neta & 12.14 \% & 11.64 \% \\
      \bottomrule
    \end{tabular}
    \micro{El costo de ventas gana participación sobre las ventas y reduce el peso relativo de la utilidad bruta.}
    \tiempo{0.8 minutos}
    \column{0.46\textwidth}
    \figpng[0.62\textheight]{F10_Barras_Porcentuales_Analisis_Vertical.png}
  \end{columns}
  \note{Cálculo vertical ABC.}
\end{frame}

\begin{frame}{Caso Empresa ABC: interpretación vertical}
  \begin{columns}[T,onlytextwidth]
    \column{0.46\textwidth}
    \footnotesize
    El costo de ventas pasa de 60.00 \% a 61.86 \%.

    \vspace{0.5em}
    Por ello, la utilidad bruta baja de 40.00 \% a 38.14 \%, mostrando presión sobre el margen.
    \tiempo{0.8 minutos}
    \column{0.50\textwidth}
    \figpng[0.64\textheight]{F11_Interpretacion_Vertical.png}
  \end{columns}
  \note{Interpretación vertical ABC.}
\end{frame}

\begin{frame}{Decisión derivada del análisis vertical}
  \begin{columns}[T,onlytextwidth]
    \column{0.48\textwidth}
    \begin{itemize}
      \item Revisar sostenibilidad del costo.
      \item Comparar márgenes por periodo.
      \item Evaluar estructura de financiamiento.
      \item Priorizar acciones correctivas.
    \end{itemize}
    \micro{Si una partida gana peso dentro de las ventas, la gerencia debe evaluar si reduce rentabilidad o eficiencia.}
    \tiempo{0.7 minutos}
    \column{0.48\textwidth}
    \figpng[0.62\textheight]{F14_Decisiones_Empresariales.png}
  \end{columns}
  \note{Decisión derivada vertical.}
\end{frame}
